快速定位并清除无线电干扰源，是恢复通信秩序和保障危险区域作业安全的重要环节。题目要求机器狗在同一地点测角误差固定、接收范围有限且源数未知的条件下，自主完成圆域内干扰源排查。本文从定位区域的几何性质出发，建立主动测向、路径规划与清除协同的决策方法。

\textbf{针对问题一}，本文建立\textbf{半平面交与最小包围圆模型}，把每次测向允许的方位范围表示为两个半平面的交集，使多次观测转化为统一的线性不等式约束。先用线性规划判断交集是否为空或无界，再枚举满足全部约束的边界交点，由最远顶点对计算区域直径。对于同直径圆的覆盖问题，证明其能够完全覆盖区域，当且仅当最小包围圆半径等于区域直径的一半。通过枚举两点直径圆与三点外接圆求得最小包围圆，并构造三次合法测向形成的等边三角形，其直径为\(38\,\mathrm{m}\)，最小覆盖半径却为\(21.94\,\mathrm{m}\)，故同直径圆不能覆盖。结合光学定位距离，进一步以最小包围圆半径不超过\(20\,\mathrm{m}\)作为存在保证清除位置的判据。

\textbf{针对问题二}，本文建立\textbf{基于贝叶斯更新的加权选点模型}。首次观测不仅限定方位，还要求源位于目标圆域及可接收范围内；不同可能源位置能够对应的接收半径也不同，例如源距超过\(1000\,\mathrm{m}\)后，只有较大的半径才能解释首次接收，因此观测后的位置分布未必均匀。以主观概率描述固定未知的源位置、接收半径与地点误差，经贝叶斯更新保留其观测后的关联，并使两次检测共享同一半径。对每个待选点分别预测强信号、无信号和正常示向度，构造反馈后的剩余区域，以最小包围圆半径的期望衡量平均定位精度，以半径不超过\(20\,\mathrm{m}\)的概率衡量本次完成的可能性，此时移至圆心即可清除，无需补测。将期望半径按\(20\,\mathrm{m}\)归一化，与未完成概率加权后最小化，采用分段高斯求积、网格初筛和多初值Nelder--Mead搜索求解。固定权重确定选点，改变权重汇集候选区域。原点、正东算例中，提高完成概率权重使优选点趋近首次点，同时增大期望半径；等权下优选点为\((813.8,\pm521.0)\,\mathrm{m}\)，期望半径为\(25.67\,\mathrm{m}\)，完成概率为\(37.34\%\)。

\textbf{针对问题三}，本文建立\textbf{基于可行区域的搜索定位清除协同模型}，将未覆盖区域的扫描和已发现源的处理作为两类任务，综合估计移动、测量与换频时间，联合安排执行顺序。采用多起点近邻与路径片段反转优化路线，每执行一项任务就根据新观测更新区域并重新规划；沿途补测兼顾多个目标，对较小区域比较继续测向与提前试清的预计耗时。逐频道检查实际无信号观测的排除圆是否完整覆盖目标圆域，结合已发现源全清判定结束。三次正式测试分别清除15、12、10个源，平均定位清除时间分别为\(195.399\)、\(276.748\)、\(248.719\,\mathrm{s}\)。两批50场、30场官方演练均全部清除，场均每源耗时分别为\(231.71\,\mathrm{s}\)和\(226.34\,\mathrm{s}\)；前者较独立演练的文献适配基线均值降低\(30.54\%\)，主要收益来自减少移动。

\textbf{针对问题四}，本文建立\textbf{朝向鲁棒搜索与联合可行域模型}，将无信号可能来自超距或背向的特点纳入源位置、接收半径和朝向的联合约束。搜索时，要求每个可能源位置位于一组站点的凸包内，且到这些站点均不超过最小接收半径，从而保证任意朝向下至少有一站可接收。由此构造22站布局，并用4988个单元校核连续圆域覆盖。定位时，利用有信号观测限制半径下界，将无法用超距解释的无信号转为背向约束；仅在不存在共同相容朝向时排除相应位置块。采用分块二分和角区间求交裁剪区域，通过有限假设估计接收概率与测向收益，再用多初始路线和2-opt改进协调行动。所选策略在30场官方演练中全部清除390个源，场均每源耗时为\(458.37\,\mathrm{s}\)，未出现光学清除失败。
